\chapter{Benchmarks de PQC no ESP32}

\section{xTaskCreatePinnedToCore}
O método \texttt{loop\(\)} padrão do Arduino limita a stack de execução para apenas 8Kb. Para isso, podemos utilizar a função xTaskCreatePinnedToCore, que nos permite criar uma thread independente e fixá-la no Core 0 (focado em Wifi, bluetooth e Sistema) ou no Core 1 (onde o sistema arduino roda por padrão). Existe também a função xTaskCreate, que deixa a decisão do core utilizado para o FreeRTOS (tskNO_AFFINITY). Os principais parâmetros dessa função a serem vistos aqui são:
\begin{itemize}
	\item usStackDepth (Tamanho da stack alocada em bytes)
	\item pvParameters (Parâmetros passados para a tarefa)
	\item uxPriority (Prioridade no agendador)
	\item pxCreatedTask (Handler para controle)
	\item xCoreId (Core 0 ou 1)
\end{itemize}

Uma alternativa para criar a task seria definir globalmente o tamanho da stack (SET_LOOP_TASK_STACK_SIZE) e colocar o código dentro do loop. Afinidade com o core 1. Precisaria criar um delay em loop ao final do código para evitar o Watchdog Timer Reset.

\subsection{usStackDepth}
Inicio todos os benchmarks com 10240 e pretendo ir reduzindo de acordo com a quantidade de memória disponível durante o teste. Se um algoritmo de segurança demorar 1,8 segundos (como o sphincs-sha2-256f-simple) dentro do loop() sem liberar tempo de processador, o Task Watchdog Timer (TWDT) do ESP32 pode interpretar que a aplicação travou e reiniciar o chip no meio do teste.
\subsection{pvParameters}
\subsection{uxPriority}
\subsection{pxCreatedTask}
\subsection{xCoreId}
